\documentclass{article}
\usepackage[utf8]{inputenc}
\usepackage{amsmath,amssymb}
\usepackage[most]{tcolorbox}
\usepackage{geometry}
\geometry{margin=2.5cm}

\newcommand{\step}[1]{\par\noindent\hangindent=3.5em\hangafter=1%
  \makebox[3.5em][l]{\textbf{#1.}}}
\newcommand{\steptag}[1]{\hfill{\small\textsf{[#1]}}}

\newtcolorbox{proofbox}[1]{
  colback=gray!5, colframe=gray!60,
  fonttitle=\bfseries, title={#1},
  breakable, enhanced,
  left=4pt, right=4pt, top=4pt, bottom=4pt,
  before skip=10pt, after skip=10pt
}

\begin{document}

\begin{proofbox}{There are universal e\_it > 0, K\_disp < infinity, and K\_fl...}
\step{1} There are universal e\_it > 0, K\_disp < infinity, and K\_floor < infinity such that, if B is a finite-dimensional C*-algebra, A is an extended epsilon-C*-algebra, and v:B->A is an extended d-inclusion with d+epsilon <= e\_it, then one dagger-preserving v\_tilde, with v\_tilde\_n = I\_n tensor v\_tilde, satisfies sup\_n ||v\_tilde\_n - v\_n|| <= K\_disp*d and is an extended K\_floor*epsilon-inclusion; for epsilon > 0 it is reached after finitely many correction steps, and for epsilon = 0 it is their operator-norm limit. \steptag{validated} \\[6pt]
\step{1.1} Let K:=K\_step>=1 and e:=e\_step>0 be the universal witnesses supplied by the registered external lem-maincb-improvement-one-step, and set e\_it:=min\{e,1/(4K)\}, K\_disp:=4K, and K\_floor:=2K. For data in the root with d+epsilon<=e\_it, put v\_0:=v and d\_0:=d. The external is applicable at j=0 and whenever d\_j<=d\_0, since d\_j+epsilon<=d\_0+epsilon<=e\_it<=e; it supplies a dagger-preserving v\_\{j+1\}, with v\_\{j+1,n\}=I\_n tensor v\_\{j+1\}, sup\_n||v\_\{j+1,n\}-v\_\{j,n\}||<=K d\_j, and an extended d\_\{j+1\}-inclusion satisfying d\_\{j+1\}<=K(d\_j\textasciicircum{}2+epsilon). If d\_j>2K epsilon, then K d\_j\textasciicircum{}2<=d\_j/4 because d\_j<=d\_0<=1/(4K), while K epsilon<d\_j/2, so d\_\{j+1\}<=(3/4)d\_j<=d\_0. Thus the correction remains legal while above the floor. If epsilon>0 and d\_0<=2K epsilon, perform exactly one correction; then d\_1<=K d\_0\textasciicircum{}2+K epsilon<=d\_0/4+K epsilon<=3K epsilon/2<=2K epsilon. \steptag{validated} \\[4pt]
\step{1.2} Assume epsilon>0. If d\_0<=2K epsilon, the single correction specified in 1.1 stops with d\_1<=2K epsilon. Otherwise repeatedly invoke lem-maincb-improvement-one-step while d\_j>2K epsilon and stop at the first N>=1 with d\_N<=2K epsilon. This N is finite because before stopping the estimate in 1.1 gives d\_j<=(3/4)\textasciicircum{}j d\_0, and a positive geometric sequence eventually falls below 2K epsilon. Every correction output is dagger-preserving, so v\_N is dagger-preserving; it is an extended d\_N-inclusion with d\_N<=K\_floor epsilon. Unpacking def-extended-delta-inclusion, each defect and norm inequality with parameter d\_N remains valid after replacing d\_N by the larger K\_floor epsilon, so v\_N is an extended K\_floor*epsilon-inclusion, with v\_\{N,n\}=I\_n tensor v\_N. \steptag{validated} \\[4pt]
\step{1.3} For the epsilon>0 construction, the displacement telescopes. In the one-step floor case it is at most K d\_0<=4K d\_0. Otherwise every preterminal defect obeys d\_j<=(3/4)\textasciicircum{}j d\_0, so sup\_n||v\_\{N,n\}-v\_\{0,n\}||<=sum\_\{j=0\}\textasciicircum{}\{N-1\}sup\_n||v\_\{j+1,n\}-v\_\{j,n\}||<=K sum\_\{j=0\}\textasciicircum{}\{N-1\}d\_j<=K d\_0 sum\_\{j>=0\}(3/4)\textasciicircum{}j=4K d\_0=K\_disp*d. \steptag{validated} \\[6pt]
\step{1.3.1} Using validated nodes 1.1 and 1.2, set K:=K\_step and K\_disp:=4K, with d\_0=d. They provide the finite epsilon>0 stopping construction v\_0,...,v\_N and, for each correction index 0<=j<N, the bounds sup\_n||v\_\{j+1,n\}-v\_\{j,n\}||<=K d\_j. If d\_0<=2K epsilon then N=1, hence sup\_n||v\_\{N,n\}-v\_\{0,n\}||<=K d\_0<=K\_disp d. If d\_0>2K epsilon, then every j<N is preterminal (d\_j>2K epsilon), so 1.1 gives d\_\{j+1\}<=(3/4)d\_j; induction yields d\_j<=(3/4)\textasciicircum{}j d\_0 for 0<=j<N. For every n, the triangle inequality gives ||v\_\{N,n\}-v\_\{0,n\}||<=sum\_\{j=0\}\textasciicircum{}\{N-1\}||v\_\{j+1,n\}-v\_\{j,n\}||. Taking sup\_n and using sup\_n(sum\_j a\_\{j,n\})<=sum\_j sup\_n a\_\{j,n\}, followed by the one-step displacement bounds, yields sup\_n||v\_\{N,n\}-v\_\{0,n\}||<=K sum\_\{j=0\}\textasciicircum{}\{N-1\}d\_j<=K d\_0 sum\_\{j=0\}\textasciicircum{}\{infinity\}(3/4)\textasciicircum{}j=4K d\_0=K\_disp d. Thus both stopping cases prove the displacement asserted in 1.3. \steptag{validated} \\[4pt]
\step{1.4} Assume epsilon=0 and apply lem-maincb-improvement-one-step indefinitely. The construction is legal by 1.1, and d\_\{j+1\}<=K d\_j\textasciicircum{}2<=d\_j/4, so d\_j<=4\textasciicircum{}\{-j\}d. For k>j, sup\_n||v\_\{k,n\}-v\_\{j,n\}||<=K sum\_\{r=j\}\textasciicircum{}\{k-1\}d\_r<=K d 4\textasciicircum{}\{-j\}/(1-1/4), which tends to zero uniformly in n. Completeness of A at level one and finite dimensionality of B give an operator-norm limit linear map v\_tilde:B->A. For each n, its amplification I\_n tensor v\_tilde is the operator-norm limit of v\_\{j,n\}=I\_n tensor v\_j, and the same tail estimate gives sup\_n||v\_tilde\_n-v\_n||<=K sum\_\{j>=0\}d\_j<=4K d/3<=K\_disp*d. \steptag{validated} \\[4pt]
\step{1.5} The epsilon=0 limit from 1.4 is dagger-preserving and is an extended 0-inclusion. Indeed every v\_j for j>=1 is dagger-preserving by lem-maincb-improvement-one-step, and the involution is isometric by def-operator-space, so dagger preservation passes to the operator-norm limit. At each fixed amplification, the two-sided (1+-d\_j) bounds in def-extended-delta-inclusion pass to the limit as d\_j tends to zero. Its d\_j-homomorphism defining relations also pass to the limit: when epsilon=0, def-extended-epsilon-cstar-algebra makes every target amplification a 0-C*-algebra, whose product is continuous with ||XY||<=||X||||Y||, and the unit, multiplication, linearity, and dagger identities therefore have zero limiting defect. Thus each v\_tilde\_n=I\_n tensor v\_tilde is a 0-homomorphism with exact two-sided norm bounds, so v\_tilde is an extended 0-inclusion, equivalently an extended K\_floor*epsilon-inclusion. \steptag{validated} \\[6pt]
\step{1.5.1} Assume epsilon=0. The already validated convergence step 1.4 supplies maps v\_j:B->A with extended defects d\_j tending to 0, a linear operator-norm limit v\_tilde:B->A, and, for every n, v\_\{j,n\}=I\_n tensor v\_j converging in operator norm to v\_tilde\_n=I\_n tensor v\_tilde (indeed uniformly in n), together with the stated displacement bound. These are precisely the existence and convergence premises used in the limit passages of node 1.5. \steptag{validated} \\[4pt]
\step{1.6} The constants e\_it, K\_disp, and K\_floor chosen in 1.1 are respectively positive, finite, and universal because the witnesses in lem-maincb-improvement-one-step are universal. For epsilon>0, 1.2 and 1.3 give a finite correction output that is dagger-preserving, has the required amplification form, is an extended K\_floor*epsilon-inclusion, and obeys the K\_disp*d displacement bound. For epsilon=0, 1.4 and 1.5 give the operator-norm limit with the same amplification identity and displacement bound, dagger preservation, and an extended 0=K\_floor*epsilon inclusion. The two cases are exhaustive and establish every assertion of the root contract. \steptag{validated} \\[6pt]
\step{1.6.1} Dependency bridge: require nodes 1.1, 1.2, 1.3, 1.4, and 1.5 to be validated. Then 1.1 supplies universal positive/finite constants e\_it, K\_disp, K\_floor and the legal correction setup. In the exhaustive case epsilon>0, 1.2 supplies a finite dagger-preserving amplified output that is an extended K\_floor*epsilon-inclusion, while 1.3 supplies its displacement bound sup\_n ||v\_\{N,n\}-v\_n|| <= K\_disp*d. In the case epsilon=0, 1.4 supplies the operator-norm limit, its amplification identity, and displacement bound, while 1.5 supplies dagger preservation and the extended 0=K\_floor*epsilon-inclusion property. Hence, only after all five dependencies are validated, these conclusions jointly imply every clause of the root contract. \steptag{validated} \\[4pt]
\end{proofbox}

\end{document}
